\section{What We Built \& What We Found}

% ---------------------------------------------------------------
\begin{frame}{Correctness \& Throughput --- Making A8 Run At All}
  \textbf{Jun 5: Automatic-8 wrote no output.} The largest instance ($10\,466$ nodes
  in a tight $342{\times}294$ canvas) has vertices sitting on foreign edges, an
  invalid layout the solver must repair first. The old $O(n\cdot m)$ repair never
  converged, so the program exited without a file.
  \medskip

  Two fixes, both shipped in the final solver:
  \begin{itemize}
    \item \textbf{Grid-accelerated repair} --- A vertex grid shares the edge-grid
          geometry, so each vertex only checks nearby edges. Multi-pass relocation
          with expanding radius; A8 now repairs in \emph{seconds}.
    \item \textbf{\texttt{restoreBest} skip} --- Wave restarts recomputed \emph{all}
          crossings: \textbf{27\% of A8's phase-1 budget}. Skip it when the state
          already matches the best $\Rightarrow$ $342$k $\to$ $1.01$M moves
          (\textbf{$3\times$}).
  \end{itemize}
  \medskip
  \footnotesize A faster per-move overlap check and an LNS time-guard round out the
  throughput work.
\end{frame}

% ---------------------------------------------------------------
\begin{frame}{The Main Lever --- Stress Initialisation}
  \textbf{Reading the winners again} --- Our SA parameters already matched the winning
  paper \emph{exactly}. The real difference was their \emph{initialisation}: a
  force-directed layout (OGDF stress majorisation + FMMM).
  \medskip

  \textbf{Our method \texttt{sa-stress}} (\texttt{tools/stress\_init.py}) --- graphviz
  \texttt{sfdp}/\texttt{neato} $\to$ scale to canvas $\to$ integer grid snap $\to$
  spiral collision resolution, then the same two-phase SA on top. Same idea as OGDF,
  an off-the-shelf tool instead.
  \medskip

  \textbf{Pick the best candidate} --- Run the layout up to 8 times, estimate each
  one's crossing density from $200$k sampled edge pairs (a full count is $\sim$205M),
  hand the sparsest to SA.
\end{frame}

% ---------------------------------------------------------------
\begin{frame}{Stress Initialisation --- Why It Wins}
  \begin{center}
    \begin{tabular}{lrrr}
      \toprule
      full budget, 8 workers & A7 & A8 & A9 \\
      \midrule
      best before stress init   & 39 & 31 & 21 \\
      \textbf{with stress init}  & \textbf{28} & \textbf{7} & \textbf{10} \\
      \bottomrule
    \end{tabular}
  \end{center}
  \medskip
  \begin{itemize}
    \item \textbf{Not luck, reproducible} --- A \emph{single} 2.5-min run reaches
          $k=8\text{--}9$ on A8. The init alone starts SA at $k=24\text{--}27$, before
          a single move --- already past GD'25 winner ($15$). Best across runs: $\mathbf{7}$.
    \item \texttt{sfdp} beat \texttt{neato} everywhere; on $n>3000$ only \texttt{sfdp} runs.
    \item One exception: dense A6 resists \emph{any} init --- more on the final slide.
  \end{itemize}
\end{frame}

% ---------------------------------------------------------------
\begin{frame}{Phase 2 --- $k$-Critical Vertex Selection}
  \textbf{The observation} --- Phase 2 minimises the \emph{maximum} crossings on a
  single edge ($k$), but vertices were picked by their \emph{total} crossings ---
  effort spread everywhere except the bottleneck.
  \medskip

  \textbf{The change} --- Bias selection towards endpoints of edges within a band
  $\beta=2$ of the current $k$, with quadratic emphasis on proximity to $k$.
  \begin{itemize}
    \item \textbf{A/B (A7, 2\,min, same seeds):} $k=55/61$ with vs.\ $77/75$ without.
    \item \textbf{Hidden scaling bug, later fixed:} a fixed base mass made $\sim$97\%
          of A8's picks effectively uniform; scaling it to the critical mass restores
          $\sim$75\% critical picks at any $n$.
  \end{itemize}
\end{frame}

% ---------------------------------------------------------------
